\documentclass{situationist-preprint}

\title{Why It Is Now}
\msid{TSM-18-a}
\msfamily{TSM-18}
\msdiscipline{Time, Space and Modality}
\mstype{Principal manuscript}
\mskeywords{now; indexicals; character; self-location; A-theory; B-theory}

\begin{document}
\maketitle

\begin{abstract}
Why is it now now, rather than at some earlier or later time? The question can seem to demand a deep answer: an explanation of why the present is located where it is in the history of the universe. We argue that it has a trivial answer and no deep one, because \enquote{now} is an indexical. We distinguish four readings of the question. Read as asking why the time of the question is the time of the question, it is answered by the character of \enquote{now}. Read as asking why this question is asked at this date, it is an ordinary causal question about the questioner. Read as asking why one time rather than another is objectively present, it has a false presupposition. Read as asking why I find myself at this time, it asks for a fact about the world that would single out one center, and there is none. No reading leaves a deep question unanswered, and there is no standpoint-free fact of what time it is now.
\end{abstract}

\section{The problem}

On the moving-spotlight A-theory, one time is lit by presentness and the light moves; it is then a real question why the spotlight is here, and another how anyone could know that it is \citep{Bourne2002,BraddonMitchell2004}. On the B-theory, \enquote{now} behaves like \enquote{here} and \enquote{I}: it refers, in each context of use, to the time of that context \citep{Smart1963,Kaplan1989}. The question \enquote{why is it now now?} then resembles \enquote{why am I here?} or \enquote{why am I this person?}, which \citet{Nagel1986} took to express a genuine puzzle about the relation between the objective and the subjective standpoints. We argue that, once its readings are separated, the question has a trivial answer and no deep one. The argument relies on TSM-15, on which the present has no metaphysical privilege, and TSM-16, on which there is no frame-independent present.

\section{Framework}

\begin{definition}[Character of \enquote{now}]\label{d:now}
In any context $c$, \enquote{now} refers to the time of $c$. A sentence containing \enquote{now} expresses, in $c$, a content about the time of $c$.
\end{definition}

\begin{definition}[Readings]\label{d:read}
The question \enquote{why is it now now?}, asked in a context $c$ whose time is $t_0$, admits four readings: (R1) why the time of $c$ is the time of $c$; (R2) why the question is asked at $t_0$; (R3) why $t_0$, rather than another time, is objectively present; (R4) why I, the questioner, am located at $t_0$ rather than at another time.
\end{definition}

\begin{principle}[No privileged present]\label{p:nopriv}
No time is objectively present: presentness is a relation of a time to what occurs at it (TSM-15), and there is no frame-independent present (TSM-16).
\end{principle}

Definition~\ref{d:now} is Kaplan's account of the character of \enquote{now} \citep{Kaplan1989}. We know of no reading of the question in the literature that is not one of (R1)--(R4) or a combination of them.

\section{The argument}

\begin{proposition}\label{p:r1}
On reading (R1), the answer is trivial: \enquote{it is now the time it now is} is true in every context.
\end{proposition}

\begin{proof}
By Definition~\ref{d:now}, in any context $c$ the sentence says that the time of $c$ is the time of $c$. Like \enquote{I am here now}, it is true in every context of use \citep{Kaplan1989}, and its truth is explained by the character of the words alone.
\end{proof}

\begin{proposition}\label{p:r2}
On reading (R2), the question is an ordinary causal question about the questioner, not a question about time.
\end{proposition}

\begin{proof}
That the question is asked at $t_0$ is a fact about a token: a person thinks a thought at a date. It is explained, as far as it is, by the causal history that brought that person to exist and think at $t_0$, as one explains why a letter was written in a given year, and the explanation says nothing about $t_0$ that it would not say about any time at which questions are asked.
\end{proof}

\begin{proposition}\label{p:r3}
On reading (R3), the question has a false presupposition.
\end{proposition}

\begin{proof}
Reading (R3) presupposes that some time is objectively present. By Principle~\ref{p:nopriv}, none is.
\end{proof}

\begin{proposition}\label{p:r4}
On reading (R4), the question asks for a fact about the world that singles out one center, and there is none.
\end{proposition}

\begin{proof}
A self-locating thought has a centered content, a set of centers each consisting of a world, a subject and a time \citep{Lewis1979}. The uncentered world contains me at many times; on the perdurantism of PID-11 it contains my temporal parts, each at its own time, and each thinks truly \enquote{I am at this time}. To ask why I am at $t_0$ is to ask why the center of this thought is $\langle w, \text{me}, t_0\rangle$ rather than another. That this thought occurs at $t_0$ is the trivial answer of Proposition~\ref{p:r1}. A further answer would have to cite a fact of the uncentered world that marks $t_0$ as the time at which I \emph{really} am, and such a fact would be the objective presentness that Principle~\ref{p:nopriv} denies.
\end{proof}

\begin{theorem}\label{t:main}
The question why it is now now has a trivial answer and no deep one. There is no fact, independent of a standpoint, of what time it is now.
\end{theorem}

\begin{proof}
By Definition~\ref{d:read}, the question has readings (R1)--(R4). By Propositions~\ref{p:r1}, \ref{p:r2}, \ref{p:r3} and~\ref{p:r4}, these have, respectively, a trivial answer, a causal answer that is not about time, a false presupposition, and no answer beyond the trivial one. A standpoint-free fact of what time it is now would answer (R3) or (R4) non-trivially, so there is none.
\end{proof}

The second part of Theorem~\ref{t:main} is what REL-19-a uses: a timeless being that does not know what time it is now is ignorant of no standpoint-free fact, because there is none to be ignorant of.

\section{Objections and replies}

\begin{objection}[The residue]\label{o:residue}
Even when the readings are separated, something remains: the vertigo of noticing that, of all the times there are, this one is mine \citep{Nagel1986}.
\end{objection}

\begin{reply}
What remains is the cognitive role of indexical thought. Thinking of a time as \enquote{now} differs from thinking of it by date, as thinking of oneself as \enquote{I} differs from thinking of oneself by name, and the difference explains action and emotion \citep{Perry1979}. It is a difference between ways of thinking of the same time, not a further fact about the time.
\end{reply}

\begin{objection}[Self-locating evidence]\label{o:evid}
Anthropic and Doomsday reasoning treat the fact that I find myself at this time as evidence \citep{Bostrom2002}. If being now were trivial, it could carry no information.
\end{objection}

\begin{reply}
Self-locating evidence is real, and its treatment is substantive, as the residual disagreement noted in EPI-18 shows. But it bears on hypotheses about the world, such as how many observers there are and when, through the probability of my center under each; it neither answers nor needs a question about which time is objectively present.
\end{reply}

\begin{objection}[The A-theory does better]\label{o:atheory}
On the A-theory the question has a deep answer: it is now now because presentness is located here.
\end{objection}

\begin{reply}
The A-theorist's answer relocates the question rather than answering it: why is presentness located at $t_0$? If the answer is that presentness moves and has reached $t_0$, the question recurs for the moving spotlight, and inhabitants of other times, thinking the same thought with the same evidence, would be mistaken; the A-theorist then cannot say how we know that we are not among them \citep{BraddonMitchell2004}. The B-theory makes every such thought true and leaves nothing to know.
\end{reply}

\section{Conclusion}

\enquote{Now} is an indexical, and the question why it is now now dissolves into a trivial truth of its character, an ordinary causal question about the questioner, a false presupposition of objective presentness, and a request for a standpoint-free fact that does not exist. It is now now because \enquote{now} refers, in each context, to the time of that context.

\begin{thebibliography}{99}\small
\bibitem[Bostrom(2002)]{Bostrom2002} Bostrom, N. (2002). \emph{Anthropic Bias: Observation Selection Effects in Science and Philosophy}. Routledge.
\bibitem[Bourne(2002)]{Bourne2002} Bourne, C. (2002). When am I? A tense time for some tense theorists? \emph{Australasian Journal of Philosophy}, 80(3), 359--371.
\bibitem[Braddon-Mitchell(2004)]{BraddonMitchell2004} Braddon-Mitchell, D. (2004). How do we know it is now now? \emph{Analysis}, 64(3), 199--203.
\bibitem[Kaplan(1989)]{Kaplan1989} Kaplan, D. (1989). Demonstratives. In J. Almog, J. Perry and H. Wettstein (eds.), \emph{Themes from Kaplan}, 481--563. Oxford University Press.
\bibitem[Lewis(1979)]{Lewis1979} Lewis, D. (1979). Attitudes de dicto and de se. \emph{The Philosophical Review}, 88(4), 513--543.
\bibitem[Nagel(1986)]{Nagel1986} Nagel, T. (1986). \emph{The View from Nowhere}. Oxford University Press.
\bibitem[Perry(1979)]{Perry1979} Perry, J. (1979). The problem of the essential indexical. \emph{No{\^u}s}, 13(1), 3--21.
\bibitem[Smart(1963)]{Smart1963} Smart, J.~J.~C. (1963). \emph{Philosophy and Scientific Realism}. Routledge \& Kegan Paul.
\end{thebibliography}

\end{document}
